\scanpage{179}
\mspage{493}

s c e n c e (*) d e \ l ' a u t r e \ t e r m e, constitue un état de déséquilibre grave, et à la limite (quand ``l'absence'' en question est quasiment complète, et prolongée) un état qui mène à la destruction de la chose (ou de l'être\unsure{)|:} en laquelle ce déséquilibre a lieu, voire même de lui et de son entourage.

Ainsi, un état d' a c t i o n ininterrompue, qui n'alterne avec des périodes suffisantes d' i n a c t i o n , de repos, mène à l'épuisement, la maladie et (à la limite) à la mort - chose d'ailleurs qui a été des plus actuelles dernièrement, pour moi! (**) Mais inversement, un état d'inaction excessif mène à un affaiblissement et à une sclérose des capacités et fonctions du corps ou du psychisme (suivant les cas), et à la limite, à la destruction. Dans le cas de mon ``incident-maladie'' d'ailleurs, j'ai un exemple simultané des d e u \unsure{x} déséquilibres: action excessive de l'esprit, inaction du corps (et un repos suffisant ni pour l'un ni pour l'autre...).

Cette ``explication'', dans ce cas d'espèce, de la ``philosophie'' équilibre-déséquilibre du yin et du yang, reste superficielle, en ce sens qu'elle ne touche pas à un parti-pris culturel invétéré, valorisant le terme\unsure{.} yang, l'action, en l' o p p o s a n t au terme yin\handwritten{,}l'inaction. Celle-ci est ressentie comme une chose ``négative'', pas productive ni intéressante à aucun point de vue, admise à la rigueur comme un pis-aller, lequel s'impose hélas même à la meilleure volonté du monde, puisqu'il faut quand-même se reposer de temps en temps pour pouvoir continuer à s'investir dans l'action (sous peine justement, comme je viens de l'expliquer, de surmenage et Dieu sait quoi encore...). En somme, l'inaction est vue comme l'humble servante de l'action, indispensable hélas mais à part çà indigne d'attention ni d'estime.

\noindent\rule{.22\linewidth}{.4pt}\par
(*) (16 octobre) En fait, cette ``absence'' n'est jamais totale il me semble - en aucune chose, le yin ni le yang n'est présent à l'état pur, sans la présence simultanée de son complémentaire, si infime soit-elle. Le ``déséquilibre'' dont je parle est donc caractérisé, non par l'absence totale d'un des deux termes complémentaires (chose jamais réalisée), mais par un état de f a i b l e s s e excessive de ce terme. Un autre type de déséquilibre, ou de morbidité, se présente quand l'un e t l'autre terme sont ``absents'', ou plus précisément, sont présents mais très faiblement. Ainsi, dans le cas du couple ``action-inaction'', un état d' a g i t a t i o n , qui n'``agit'' pas à proprement parler (si ce n'es\unsure{t} pour se perpétuer lui-même, pour maintenir une confusion), tout en dispersant de l'énergie, peut être considéré sans doute comme un tel déséquilibre ``par défaut'' (du yin e t du yang).

(**) Voir à ce sujet les deux premières notes (n$^\circ$ 98, 99) du Cortège XI, ``Le défunt (toujours pas décédé...)''.

% archive sheet number: 178
